\section{引言}
\label{sec:intro}

随着信息技术的持续进步，图像已经成为了人们信息存储和传输的重要载体，其广泛应用于互联网、医疗影像、卫星遥感国防安全等重要领域。[来一篇综述]数字图像在传输和存储过程中容易受到各种攻击，如窃取、篡改和伪造等，导致图像信息的泄露和损坏，造成不可估计的损失。因此，图像加密技术作为保护图像信息安全的重要手段，受到了广泛关注。[来两篇参考文献]

这一段写CLM混沌系统在图像加密中的应用现状
耦合映射格子（Coupled Map Lattice, CML）是Kaneko提出的经典时空混沌系统\cite{Kaneko1989Spatiotemporal}，它以离散时间、离散空间、连续状态的格点阵列为载体，通过近邻扩散耦合将局部映射的混沌行为在空间维度上传播，从而在保持迭代形式简洁的同时获得高维相空间与多个正 Lyapunov 指数 [一篇文献]围绕标准 CML 存在的耦合方式单一、格点间同步倾向明显、参数空间中混沌区域不连续等缺陷，后续研究者提出了多种改进方案。Zhang 与 Wang 提出混合线性–非线性耦合 Logistic 映像格子（MLNCML），通过在耦合项中同时引入线性与非线性分量显著扩大了混沌参数区间并抑制了空间同步 [7,2014]，并据此构建了对称图像加密算法 [8,2014]；随后又以非邻耦合（non-adjacent coupling）打破空间局域性，使扰动能够跨格点快速扩散 [9,2015]。在此基础上，引入耦合参数 q 的广义格子 [10,2017]、动态耦合系数模型（DCCML）[11,2018]、非邻动态耦合模型 [12,2019]、交叉耦合映像格子（CCML）[13,2020; 14,2020]、多重动态耦合格子 [15,2020]、分段多动态耦合系数格子 [16,2021] 以及全混沌耦合映像格子 [17,2022] 相继被提出，其共同思路是令耦合强度本身由混沌轨道驱动、随时间与空间位置变化，从而消除固定耦合参数带来的可预测性并提升 Kolmogorov–Sinai 熵与排列熵等复杂度指标。与此并行的另一改进方向是替换或复合局部映射，Logistic–Chebyshev 动态耦合格子 [18,2021]、混合多混沌耦合格子 [19,2020]、LSS 型耦合格子 [20,2024]，以及将模型由一维推广至二维的 2D-MLNCML 与 2DNLCML 系统 [21,2018; 22,2020]，后者在结构上与图像的行列组织进一步对齐。这些改进模型解决了传统模型局部耦 合范围有限的问题,但忽视了参数设置问题,导致混沌性质 只能在狭窄的参数范围内达到次优表现。

这一段写DNA加密的研究现状
DNA 分子具有海量存储密度、天然并行性与超低能耗等特性，在信息加密领域展现出独特的优势\cite{Yan2023ColorDNA,Afify2023DynamicDNA,Yu2024DNANetwork,Aldayel2025LLEVerified,Huang2025DNAEncoding}。zhu等\cite{zhu2023parallelDNA}结合SHA-3明文关联机制与并行DNA编解码技术，设计了一种兼具高安全性和低时间开销的快速图像加密方案，有效地降低了计算复杂度。Cun\cite{Cun2023DNARNA}等提出了基于动态DNA编码与RNA计算的四维超混沌图像加密算法，结合置换序列生成器实施行列循环移位，提升了系统的抗攻击能力。

这一段写图像加密置乱+扩散的研究现状。[一篇细粒度加密的文献]
Fridrich (1998) 提出的两步加密算法——扰乱与扩散,依然是最主流的加密方法。扰乱是图像像素的无序重排,扩散是像素值的意外变化。随着时间推移,攻击者已熟悉该两阶段加密方案,因此一些简单或经典算法安全性较低。年来,部分学者提出了比特级扰乱加密:将像素值转为二进制后操作。大量实验结果表明,独特 的比特特化扰乱方法安全性高于像素级扰乱。

[TODO]

1.提出了一种新的S-NCML时空混沌系统，在广泛的参数范围内具有更强的混沌性和更高的复杂性，能够有效地提高图像加密的安全性。

2.提出了一种分块的DNA编码和运算方法，将彩色图像分为若干个块，并对每个块进行不同的DNA编码和运算，进一步提高了系统的安全性和抗攻击能力。

3.提出来一种新的置乱扩散操作。

4.通过一系列严格的仿真实验和安全性分析,验证了  该算法在统计特征、抗攻击能力和运算效率方面的优越性,证明其能够有效地保护彩色图像的安全性和隐私性，具有广泛的应用前景。